\chapter{Improvements and Future Development}

\label{chap:Chapter8}

The implementation presented in Chapter~7 demonstrates that selected functionality from Boost.Align and Boost.Endian can be reproduced in 
Rust while maintaining interoperability with existing C++ code. The implementations provide a basis for evaluating the migration 
methodology proposed in this dissertation. However, the current implementations also identify several opportunities for improvement.

These improvements concern both the completeness of the migrated functionality and the robustness of the interfaces used to integrate the 
Rust implementations with C++. They also concern portability, error handling, testing, performance evaluation, and the reduction of assumptions 
that are specific to the current experimental environment.

The objective of this chapter is therefore not to introduce additional functionality solely for the purpose of reproducing the complete 
Boost APIs. Instead, the improvements are selected according to the objectives of the migration methodology: preserving relevant behavior, 
improving safety, maintaining interoperability, supporting portability, and providing sufficient evidence for evaluating the suitability of 
Rust for low-level automotive software.

The improvements are divided into the two selected libraries. Section \ref{sec:align-improvements} discusses improvements to the Boost.Align 
reimplementation, while Section \ref{sec:endian-improvements} discusses improvements to the Boost.Endian implementation. 
Section \ref{sec:ffi-improvements} then considers improvements applicable to the interoperability layer shared by both implementations. 
Finally, Section~\ref{sec:future-validation} presents improvements to the validation and evaluation process.

\section{Boost.Align Improvements}

\label{sec:align-improvements}

The Boost.Align implementation described in Section~7.3 provides an aligned buffer abstraction, compile-time alignment through Rust's 
\texttt{\#[repr(align(...))]} attribute, AVX-based processing, and a C-compatible interface. The implementation successfully demonstrates 
that the alignment requirements of the selected use case can be expressed using Rust's type system.

The current implementation is nevertheless intentionally narrower than the complete Boost.Align API. The following improvements would increase 
its generality and suitability for production use.

\subsection{Generalisation of Alignment Requirements}

The current \texttt{Aligned<T>} abstraction uses a fixed 32-byte alignment:

\begin{verbatim}
#[repr(align(32))]
pub struct Aligned<T>(pub T);
\end{verbatim}

This is appropriate for the AVX workload evaluated in this project because a 256-bit AVX register contains eight 32-bit floating-point values 
and requires 32-byte alignment for the aligned load and store operations used by the implementation.

However, fixing the alignment to 32 bytes limits the abstraction to one particular use case. A more general implementation should allow the 
required alignment to be selected according to the type or operation being performed.

The existing \texttt{aligned\_type!} macro partially addresses this limitation by allowing a user-defined alignment:

\begin{verbatim}
aligned_type!(TypeName, InnerType, Alignment);
\end{verbatim}

A future implementation could replace the macro-based interface with a more general abstraction based on Rust's type system and allocation 
facilities. Such an abstraction would allow the same implementation to support 8-, 16-, 32- or 64-byte alignment where required.

This would make the Rust implementation more closely applicable to the wider class of alignment requirements supported by Boost.Align, 
whose purpose is to provide portable control and inspection of memory alignment. \cite{boostalign}

\subsection{Expansion of the Boost.Align API}

The current implementation focuses on aligned storage and aligned SIMD processing rather than reproducing the complete Boost.Align API.

In particular, the implementation does not currently provide direct Rust equivalents for all of the following Boost.Align facilities:

\begin{itemize}
\item alignment calculation;
\item \texttt{align\_up};
\item \texttt{align\_down};
\item alignment inspection;
\item \texttt{is\_aligned};
\item general aligned allocation and deallocation;
\item aligned allocators; and
\item allocator adaptors.
\end{itemize}

These facilities could be added as separate Rust abstractions rather than as a direct translation of the original C++ API.

For example, alignment calculations could be expressed using integer operations, while allocation could be represented using Rust's 
\texttt{Layout} abstraction. The objective should remain to preserve the semantic guarantees of the original functionality rather than reproduce 
its C++ interface literally.

This would also allow a more direct comparison between the facilities provided by Boost.Align and their Rust equivalents.

\subsection{Removal of Unnecessary Raw Slice Construction}

The current \texttt{AlignedBuffer} implementation provides the following operations:

\begin{verbatim}
pub fn as_slice(&self) -> &[f32]
pub fn as_mut_slice(&mut self) -> &mut [f32]
\end{verbatim}

These operations construct Rust slices from raw pointers using \texttt{std\:\:slice\:\:from\_raw\_parts} and \texttt{from\_raw\_parts\_mut}. 
Consequently, they require \texttt{unsafe} code even though the underlying buffer is owned by a safe Rust abstraction.

A future implementation should investigate whether the internal representation can be reorganised so that the safe slice interface does not 
require raw-pointer reconstruction.

For example, the internal representation could use an aligned allocation whose ownership and length are represented explicitly while exposing a 
safe slice through a carefully designed abstraction. The goal would be to reduce the amount of unsafe code while preserving the alignment 
guarantee.

This would strengthen one of the central objectives of the migration strategy: minimizing and isolating unsafe Rust code.

\subsection{Improved FFI Error Handling}

The current FFI implementation uses assertions for invalid pointers and invalid buffer lengths. For example, the FFI functions verify that a 
buffer pointer is not null and that the buffer length is compatible with the AVX operation.

Assertions are useful during development because they expose violations of the expected preconditions immediately. However, they are not an 
ideal error-handling mechanism for a production C++ interface.

A future interface should replace assertions at the FFI boundary with explicit error handling. Possible approaches include returning an 
integer status code, returning a nullable pointer where appropriate, or exposing a dedicated error enumeration.

For example:

\begin{verbatim}
enum AlignError {
NullPointer,
InvalidLength,
UnsupportedArchitecture,
InvalidAlignment
}
\end{verbatim}

Such an interface would allow the C++ caller to handle invalid input without depending on Rust panic behavior.

This is particularly important for safety-oriented systems, where error handling should be deterministic and explicitly defined at component 
boundaries.

\subsection{Runtime SIMD Feature Detection}

The current SIMD implementation is explicitly compiled with:

\begin{verbatim}
#[target_feature(enable = "avx")]
\end{verbatim}

and uses AVX intrinsics such as \texttt{\_mm256\_load\_ps}, \texttt{\_mm256\_add\_ps}, and \texttt{\_mm256\_store\_ps}.

This is suitable for the experimental environment used by the project but introduces an architectural assumption: the target processor must 
support AVX.

A more portable implementation should distinguish between the aligned storage abstraction and the instruction set used to process that storage.

A future implementation could therefore provide multiple execution paths:

\begin{verbatim}
aligned buffer
|
+---- scalar implementation
|
+---- SSE implementation
|
+---- AVX implementation
|
+---- AVX2 implementation
\end{verbatim}

The appropriate implementation could then be selected according to compile-time target information or runtime CPU feature detection.

This would make the alignment component more portable across the different processor architectures and configurations encountered in 
automotive systems.

\subsection{Generalisation Beyond \texttt{f32}}

The current \texttt{AlignedBuffer} is specifically designed around \texttt{f32} values and the \texttt{F32x8} representation.

This design is appropriate for demonstrating AVX processing, but it means that the buffer abstraction itself is tied to one data type.

A future implementation could separate the generic aligned storage abstraction from the particular SIMD kernel. For example:

\begin{verbatim}
AlignedBuffer<T, A>
\end{verbatim}

could represent aligned storage independently of whether the stored values are floating-point values, integers, or user-defined structures.

The SIMD operation could then be implemented separately for each supported type and architecture.

Such a design would make the alignment component more closely resemble the general-purpose nature of Boost.Align while retaining specialized 
optimized implementations where appropriate.

\subsection{Improved Memory and Performance Evaluation}

The current benchmark demonstrates that the Rust implementation can achieve comparable performance to the Boost.Align implementation for the 
selected AVX workload. The repository reports an average iteration time of approximately 20.7~ms for Rust compared with approximately 21.7ms 
for Boost.Align.

Although this result is encouraging, the experiment evaluates a single workload and processor configuration.

A more comprehensive evaluation should vary:

\begin{itemize}
\item buffer size;
\item alignment;
\item data type;
\item CPU architecture;
\item SIMD instruction set;
\item compiler optimization level;
\item allocation frequency; and
\item number of repeated operations.
\end{itemize}

This would allow the performance conclusions to be generalized beyond the specific AVX workload used in the current implementation.

\section{Boost.Endian Improvements}

\label{sec:endian-improvements}

The Boost.Endian implementation provides the three principal usage models identified in Boost\.Endian: conversion functions, buffer types and 
arithmetic types. The implementation also provides C++ interoperability through the \texttt{cxx} crate.

The current implementation is therefore broader than the Boost.Align implementation in terms of API structure. Nevertheless, several 
improvements would increase its completeness, robustness and suitability for interoperability.

\subsection{Stabilisation of the Baseline Implementation}

The repository contains several Endian implementation variants developed during the migration process. The comparison documentation identifies 
the original \texttt{endian} crate as the Rust-native baseline, an \texttt{endianclv1} implementation using a C ABI, an \texttt{endianclv2} 
implementation using a \texttt{cxx} bridge, and an \texttt{endianv3} implementation providing the richest demonstration interface.

This experimentation is valuable because it demonstrates that FFI design is itself an important part of the migration process.

However, the baseline implementation should be consolidated before being treated as the final migrated component. In particular, the repository 
documentation records that the original Rust-native implementation still requires further stabilization.

A future version should establish one canonical Endian implementation and clearly separate experimental alternatives from the final migration 
result.

The final implementation should then be evaluated using a reproducible test and benchmark procedure.

\subsection{Expansion of Non-Standard Integer Widths}

Boost.Endian supports integer representations whose widths do not necessarily correspond to standard C++ integer types. This includes 24-, 40-, 
48- and 56-bit representations.

The current Rust implementation provides explicit 24-bit load and store functionality, with the value represented internally using a wider 
native integer type.

A natural extension would be to generalize this mechanism to all non-standard widths supported by Boost\.Endian:

\begin{verbatim}
24-bit
40-bit
48-bit
56-bit
\end{verbatim}

The representation should continue to use byte-oriented storage rather than relying on a native integer with a matching size, since Rust does 
not provide native integer types for these widths.

This would allow the Rust implementation to cover a larger portion of the Boost.Endian buffer API while retaining the portability advantages of 
explicit byte representations.

\subsection{Aligned and Unaligned Buffer Types}

The Boost.Endian buffer interface distinguishes aligned and unaligned representations. The official Boost documentation recommends unaligned 
buffers when portability is the primary concern because they avoid padding and architecture-dependent alignment requirements. cite{boostendianbuffers}

The current Rust implementation primarily represents endian-aware values through transparent wrappers around primitive types.

A future implementation could make the distinction between aligned and unaligned representations explicit.

An unaligned representation could use:

\begin{verbatim}
[u8; N]
\end{verbatim}

as its underlying storage.

An aligned representation could use an appropriately aligned wrapper where the target architecture and use case justify the additional 
constraint.

This would allow the migration to reproduce another important design choice in Boost.Endian while allowing applications to choose portability 
or potentially more efficient aligned access according to their requirements.

\subsection{Expansion of the C++ FFI Surface}

The current \texttt{cxx} bridge exposes a selected subset of the Rust implementation. In particular, the demonstrated bridge includes 
load/store and buffer functionality for commonly used integer widths.

The Rust implementation contains functionality that is not currently exposed through the C++ bridge, including the 24-bit conversion functions.

A future version should define an explicit FFI coverage matrix:

\begin{tabular}{|l|c|c|}
\hline
Feature & Boost.Endian & Rust \\
\hline
8-bit conversion  & Yes     & Planned \\
16-bit conversion & Yes     & Yes \\
32-bit conversion & Yes     & Yes \\
40-bit conversion & Planned & Planned \\
56-bit conversion & Planned & Planned \\
64-bit conversion & Yes     & Yes \\
\hline
\end{tabular}

This would make the interoperability scope explicit and prevent the C++ interface from being interpreted as a complete representation of 
the Rust implementation.

\subsection{Improved C++ Arithmetic Interface}

Rust implements the arithmetic functionality through Rust operator traits such as \texttt{Add}, \texttt{Sub}, \texttt{Mul}, \texttt{Div}, 
\texttt{BitAnd}, \texttt{BitOr} and \texttt{BitXor}. These traits provide an idiomatic Rust representation of Boost.Endian's arithmetic types.

However, operator traits cannot simply be exported through the FFI as C++ operator overloads.

The current solution therefore uses a C++ wrapper around the FFI functions.

A future implementation could improve this interface by providing a dedicated C++ compatibility layer that exposes an API closer to the 
corresponding Boost.Endian arithmetic types.

The C++ layer would remain thin, with the actual endian representation and conversion logic remaining in Rust.

This would produce the following architecture:

\begin{verbatim}
C++ compatibility type
|
v
cxx bridge
|
v
Rust EndianInt
|
v
Encoding + conversion
\end{verbatim}

Such an arrangement would allow existing C++ applications to migrate individual uses of Boost.Endian without requiring immediate changes to 
their surrounding code.

\subsection{Explicit Validation of Binary Layout}

Endian-aware data types are frequently used in binary protocols and file formats. Therefore, numerical correctness alone is insufficient to 
validate the implementation.

The tests should also verify the exact byte-level representation.

For example, a value:

$$
0x01020304
$$

stored as big endian must produce:

\begin{verbatim}
01 02 03 04
\end{verbatim}

while the corresponding little-endian representation must produce:

\begin{verbatim}
04 03 02 01
\end{verbatim}

The current C++ demonstration already exercises mixed-endian structures and checks their binary representation. This should be expanded into a 
systematic automated test suite.

The test suite should verify:

\begin{itemize}
\item byte representation;
\item structure size;
\item field offsets;
\item alignment;
\item signed values;
\item negative values for non-standard widths;
\item floating-point representations;
\item round-trip conversions; and
\item mixed-endian structures.
\end{itemize}

These tests are particularly important because a numerically correct conversion can still result in an incorrect binary layout.

\subsection{Improved FFI Error Handling}

The current Endian FFI uses Rust slices to represent buffers crossing the C++/Rust boundary.

This is considerably safer than reconstructing arbitrary Rust slices from raw pointers, but the caller is still responsible for supplying a 
buffer with the expected number of bytes.

For example, a 64-bit load requires at least eight bytes.

A future FFI layer should therefore expose fallible versions of the load and store operations. Instead of allowing an invalid slice length to 
trigger a panic, the interface could return an explicit status:

\begin{verbatim}
SUCCESS
INVALID_LENGTH
INVALID_ARGUMENT
UNSUPPORTED_TYPE
\end{verbatim}

This would provide a more deterministic interface for C++ applications and would be preferable for production-oriented integration.

\subsection{FFI Boundary Optimisation}

The repository contains a dedicated benchmark comparing a native C++ implementation with the Rust implementation accessed through the 
\texttt{cxx} bridge.

The benchmark is valuable because it evaluates not only the byte-order algorithm but also the cost of crossing the language boundary.

However, the benchmark should distinguish between two different sources of overhead:

\begin{enumerate}
\item the cost of the endian conversion itself; and
\item the cost of invoking the Rust implementation through FFI.
\end{enumerate}

The current benchmark performs a Rust load and Rust store for each element. Consequently, the number of FFI calls grows with the number of 
processed values.

A more representative systems-level implementation could batch operations so that one FFI call processes an entire buffer.

For example:

\begin{verbatim}
C++ buffer
|
| one FFI call
v
Rust process_buffer()
|
| process N elements
v
C++ buffer
\end{verbatim}

This would move the FFI boundary outside the inner processing loop and provide a more realistic model for high-throughput binary processing.

It would also allow the benchmark to distinguish the cost of the interface from the cost of the actual endian conversion.

\section{Common FFI Improvements}

\label{sec:ffi-improvements}

Although Boost.Align and Boost.Endian have different implementation characteristics, both expose similar challenges at the C++/Rust boundary.

The first improvement is to define the ownership model explicitly for every FFI object. Any pointer crossing the boundary should have a 
documented owner, lifetime and destruction mechanism.

The second improvement is to avoid allowing Rust panics to become the normal error-reporting mechanism for invalid C++ input. FFI functions 
should instead return explicit status values or use a documented error representation.

The third improvement is to minimize the amount of data exchanged through the FFI boundary. Large buffers should preferably remain owned by the 
caller while Rust operates on borrowed slices or pointers whose validity is established by the interface contract.

The fourth improvement is to keep unsafe code localized inside the FFI layer. Safe Rust abstractions should validate the relevant invariants 
before invoking the low-level operations.

This produces the general migration boundary:

\begin{verbatim}
+----------------------+       +----------------------+
|      C++ system      |       |     Rust library     |
|                      |       |                      |
| Existing application | <---> | Safe abstractions    |
| Existing data types  |  FFI  | Type-level invariants|
|                      |       | Low-level operations|
+----------------------+       +----------------------+
\                   /
\                 /
+---------------+
FFI boundary
ownership/errors/
representation
\end{verbatim}

Such a boundary is consistent with the migration objectives defined earlier in this dissertation, particularly the requirement to maintain C++ 
interoperability while minimizing and isolating unsafe Rust.

\section{Improvements to Validation}

\label{sec:future-validation}

The current validation demonstrates functional correctness, binary representation, alignment and selected performance characteristics. 
A more comprehensive evaluation should extend these tests in several dimensions.

\subsection{Cross-Architecture Testing}

Both selected libraries interact closely with processor architecture.

For Boost.Align, this concerns alignment requirements and SIMD instruction availability.

For Boost.Endian, this concerns native byte order and the relationship between the external representation and the host representation.

The implementation should therefore be evaluated on both little-endian and big-endian architectures where practical.

This would verify that the implementation does not accidentally depend on the little-endian architecture used by the development environment.

\subsection{Compiler and Toolchain Testing}

The C++ interoperability layer should also be evaluated using multiple compiler and toolchain configurations.

At minimum, the evaluation should consider:

\begin{itemize}
\item GCC;
\item Clang;
\item different supported Rust toolchains;
\item debug and release builds; and
\item different optimization levels.
\end{itemize}

This is particularly relevant because Boost libraries were historically designed to compensate for differences between C++ compilers and 
standard library implementations.

A Rust migration should therefore demonstrate that the new implementation does not simply replace one platform-specific dependency with another.

\subsection{Property-Based Testing}

The existing tests use specific representative values. These are useful for validating known cases but do not exhaustively explore the possible 
input space.

Property-based testing could provide stronger validation.

For endian conversion, an important property is:

$$
reverse(reverse(x)) = x
$$

for all supported values \(x\).

Similarly, for a supported endian representation \(E\):

$$
E^{-1}(E(x)) = x
$$

should hold for all valid values.

For aligned buffers, properties can include:

$$
address \bmod alignment = 0
$$

and:

$$
length \bmod SIMD\_width = 0
$$

for buffers passed to the SIMD implementation.

These properties can complement the existing example-based tests.

\section{Future Integration with Automotive Software}

The selected improvements also provide a path toward evaluating the migrated libraries in a more realistic automotive software environment.

For Boost.Align, future work could investigate aligned buffers in workloads involving:

\begin{itemize}
\item sensor processing;
\item SIMD-based signal processing;
\item DMA-compatible memory;
\item memory pools; and
\item performance-critical data structures.
\end{itemize}

For Boost.Endian, representative workloads could include:

\begin{itemize}
\item network protocol parsing;
\item diagnostic messages;
\item serialization;
\item binary file formats; and
\item communication between heterogeneous processors.
\end{itemize}

These scenarios would allow the migration methodology to be evaluated beyond isolated library benchmarks and closer to the types of low-level 
functionality encountered in automotive software architectures.

The Eclipse S-CORE context described earlier in this dissertation provides one potential environment in which these evaluations could eventually 
be performed.

\section{Summary}

The implementations presented in Chapter~7 demonstrate the feasibility of reimplementing selected Boost functionality in Rust. The next stage is 
therefore not necessarily to reproduce every feature of the original libraries, but to improve the implementations according to the requirements 
identified during the migration.

For Boost.Align, the principal improvements are the generalisation of alignment requirements, expansion of the API, reduction of unsafe slice 
construction, improved FFI error handling, runtime SIMD feature handling, support for additional data types, and broader performance evaluation.

For Boost.Endian, the principal improvements are stabilisation of the selected implementation variant, expansion of non-standard integer widths, 
explicit aligned and unaligned representations, broader C++ FFI coverage, improved arithmetic interoperability, stronger binary-layout 
validation, explicit FFI error handling, and reduction of per-element FFI overhead.

Across both libraries, the most important architectural improvement is the definition of a robust interoperability boundary. 
Ownership, representation, lifetime, error handling and safety invariants should be explicit at this boundary.

These improvements provide a path from the current research implementation toward a more general and production-oriented Rust library while 
preserving the central objective of the dissertation: demonstrating how low-level Boost functionality can be migrated incrementally to Rust 
without requiring an immediate replacement of the surrounding C++ system.
